\section{Verification and formalization}\label{sec:verification}
The only computer-assisted step in the proof of Theorem~\ref{thm:main} is Lemma~\ref{lem:certificate}. Everything else, including the interpretation of the certificate and the passage to arbitrary graphs, is proved in the preceding sections. Separately, the whole of Corollary~\ref{cor:allodd} has been formalized in Lean~4.

\subsection{The rational certificate}
The certificate was found by semidefinite optimization and then reconstructed over the rationals. The independent verifier enumerates the admissible graphs and flags from scratch, forms the Gram matrices from their rational factorizations, and checks all $1436$ coefficient identities~\eqref{eq:rowidentity}. Exact elimination gives $456$ positive pivots in the reduced matrices. The verifier also confirms that the multipliers and slacks are nonnegative and that $\max_\theta\alpha_\theta=383936867/10^9<1$, the bound used in Theorem~\ref{thm:stable}. Together with the averaging argument of Section~\ref{sec:certificate}, these finite checks prove Theorem~\ref{thm:polynomial}. The verifier and the certificate data are in the directory \code{certificate} of the accompanying repository \url{https://github.com/Asad-Shahab/erdos-809-lean}, and Appendix~\ref{app:data} describes the data further.

\subsection{The Lean development}
The formalization is written in Lean~4.28.0~\cite{Lean4} against mathlib~v4.28.0~\cite{Mathlib}, and is contained in the same repository; the description below refers to commit \code{7ec4aaa}. Its three main theorems are the following.
\begin{description}[leftmargin=1.5em,itemsep=3pt]
\item[\code{Erdos809.erdos_809_C7}] Theorem~\ref{thm:main}. Its proof includes the finite certificate and its interpretation, the reduction from colorings to the palette bound, and the upper bound. It is stated in \code{Erdos809/Main.lean}.
\item[\code{Erdos809.erdos_809_long_odd_cycles}] The case $k\ge4$ of Corollary~\ref{cor:allodd}, following~\cite[Sections~3--4]{BCM2026}. The mathematics of this part is due to Buci\'c, Chen, and Ma; the contribution here is its formalization. It is stated in \code{Erdos809/LongOddCycles/Main.lean}.
\item[\code{Erdos809.erdos_809}] Corollary~\ref{cor:allodd}, deduced from the two previous theorems by separating the case $k=3$ from $k\ge4$. It is stated in \code{Erdos809/OddCycles.lean}.
\end{description}

The last theorem states that for every integer $k\ge3$ and real $\varepsilon>0$ there is $N$ such that, for every integer $n\ge N$, the minimum $q$ in~\eqref{eq:fdefinition} with $H=C_{2k+1}$ and $e=\lfloor n^2/4\rfloor+1$ is attained and satisfies $(1/8-\varepsilon)n^2\le q\le(1/8+\varepsilon)n^2$. Here $N$ may depend on $k$ and $\varepsilon$. A copy of the cycle is an injective adjacency-preserving map from $C_{2k+1}$ into the host graph, so copies are simple cycles that need not be induced, and the number of colors of a coloring is the size of its image. The minimum ranges over all graphs with at least $e$ edges and all valid colorings; a separate lemma shows that it is also attained by a graph with exactly $e$ edges.

In the $C_7$ part, the certificate appears as explicit data in the directory \code{Erdos809/Certificate}, scaled to integers by the common denominator given in Appendix~\ref{app:data}. These data comprise the enumeration of admissible five-vertex graphs with their orbit representatives, the fifteen Gram blocks together with integer witnesses of their positive semidefiniteness, and the coefficient identities. All of these finite statements are checked by the Lean kernel. The weighted and asymptotic bounds are then derived from them as in Sections~\ref{sec:certificate}--\ref{sec:completion}, with the regularity lemma and the triangle removal lemma taken from mathlib. The final theorem depends only on the standard axioms \code{propext}, \code{Classical.choice}, and \code{Quot.sound}; in particular, no step relies on compiled code or on an external computation.

The formalization differs from the exposition above in two inessential respects. It extracts a triangular matching from each exact cover, and its two-clique construction uses clique sizes differing by an $\varepsilon$-dependent linear amount. These variants give the same palette lower bound and the same asymptotic upper bound as the arguments presented here.
